% Source: ne630_problems/lesson_28, HEAD 4f375f82; see homework/README.md.
% Current homework/markdown/HW28.md is authoritative; refreshed 2026-11-02.
% FNRP 5.9 attribution retained privately here; Problem 3 is no longer a bonus.

\noindent{\bf Problem 1}

Find the periods for reactors fueled by uranium-235, plutonium-239, and
uranium-233 if

\begin{itemize}
\setlength{\itemsep}{0pt}\setlength{\parskip}{0pt}
\item
  \textbf{(a)} One cent of reactivity is added to the critical systems.
\item
  \textbf{(b)} One cent of reactivity is withdrawn from the critical
  systems.
\end{itemize}

\begin{adaptedsolution}
One cent is $\rho=\pm0.01\beta$. As in the source solution, use the
small-reactivity approximation and neglect prompt generation time compared
with the delayed-neutron contribution:
\[
T\simeq\frac{\Lambda+\sum_i\beta_i/\lambda_i}{\rho}
\simeq\frac{\beta}{\lambda_{\rm eff}\rho},\qquad
\lambda_{\rm eff}=\frac{\sum_i\beta_i}{\sum_i\beta_i/\lambda_i}.
\]
Using the Table 5.1 data recorded in \texttt{HW28\_sol.ipynb} gives
\begin{center}
\begin{tabular}{lrrr}
\toprule
Fuel & $\lambda_{\rm eff}$ ($\mathrm{s^{-1}}$) & $T$ for $+1$ cent (s)& $T$ for $-1$ cent (s)\\
\midrule
U-235 & 0.0766409 & 1304.79 & $-1304.79$\\
Pu-239 & 0.0636875 & 1570.17 & $-1570.17$\\
U-233 & 0.0494377 & 2022.75 & $-2022.75$\\
\bottomrule
\end{tabular}
\end{center}
These are approximately $\pm21.8$, $\pm26.2$, and $\pm33.7$ minutes,
respectively. A negative period denotes decay, not a negative elapsed time.
The prompt does not specify $\Lambda$; therefore these are the source's
small-reactivity estimates, not exact roots of a fully specified inhour equation.
\end{adaptedsolution}

\newpage

\noindent{\bf Problem 2}

Consider a critical reactor with \(n(0)=10^6~\mathrm{cm^{-3}}\) subject
to a step insertion of reactivity \(\rho_0\) at \(t=0\). Using the
kinetics equations with a single precursor group, i.e.,

\[
\begin{aligned}
\frac{dn}{dt}&=\frac{\rho(t)-\beta}{\Lambda}n(t)+\lambda c(t),\\
\frac{dc}{dt}&=-\lambda c(t)+\frac{\beta}{\Lambda}n(t),
\end{aligned}
\]

compute \(n(t)\) and \(c(t)\) at \(t=0\), \(t=0.1\), and
\(t=10~\mathrm{s}\) assuming \(\beta=0.0065\),
\(\lambda=0.08~\mathrm{s^{-1}}\), \(\Lambda=5\times10^{-5}~\mathrm{s}\),
and

\begin{itemize}
\setlength{\itemsep}{0pt}\setlength{\parskip}{0pt}
\item
  \textbf{(a)} \(\rho_0=0.2\beta\)
\item
  \textbf{(b)} \(\rho_0=-0.2\beta\)
\end{itemize}

Make sure to clearly indicate intermediate results, e.g., the time
constants, coefficients, etc.

\begin{adaptedsolution}
The neutron equation must include the factor $n(t)$ multiplying
$(\rho-\beta)/\Lambda$; its omission in the source prompt is a typographical
error, as confirmed by the source derivation and notebook.
Initially $c_0=\beta n_0/(\lambda\Lambda)=1.625\times10^9~\mathrm{cm^{-3}}$.
For $t\ge0$, the characteristic equation is
\[
\omega^2+B\omega+C=0,\qquad
B=\lambda-\frac{\rho_0-\beta}{\Lambda},\qquad
C=-\frac{\lambda\rho_0}{\Lambda}.
\]
Let $\omega_1>\omega_2$ be its roots. The source solution gives
\[
\begin{aligned}
n(t)&=a_1e^{\omega_1t}+a_2e^{\omega_2t},&
c(t)&=b_1e^{\omega_1t}+b_2e^{\omega_2t},\\
a_1&=\frac{n_0(\rho_0/\Lambda-\omega_2)}{\omega_1-\omega_2},&
a_2&=n_0-a_1,\\
b_1&=\frac{-c_0\omega_2}{\omega_1-\omega_2},&b_2&=c_0-b_1.
\end{aligned}
\]
Here $\omega_i$ are exponential rates in $\mathrm{s^{-1}}$ (the source
calls them time constants); their reciprocals are modal periods.
\begin{center}
\begin{tabular}{lrr}
\toprule
Quantity & $\rho_0=+0.2\beta$ & $\rho_0=-0.2\beta$\\
\midrule
$B$ ($\mathrm{s^{-1}}$) & 104.08 & 156.08\\
$C$ ($\mathrm{s^{-2}}$) & $-2.08$ & 2.08\\
$\omega_1$ ($\mathrm{s^{-1}}$) & 0.0199807914 & $-0.0133276373$\\
$\omega_2$ ($\mathrm{s^{-1}}$) & $-104.099981$ & $-156.066672$\\
$a_1$ ($\mathrm{cm^{-3}}$) & 1249520.06 & 833475.71\\
$a_2$ ($\mathrm{cm^{-3}}$) & $-249520.06$ & 166524.29\\
$b_1$ ($\mathrm{cm^{-3}}$) & $1.62468816\times10^9$ & $1.62513878\times10^9$\\
$b_2$ ($\mathrm{cm^{-3}}$) & 311840.17 & $-138782.10$\\
\bottomrule
\end{tabular}
\end{center}
\begin{center}
\begin{tabular}{crrrr}
\toprule
&\multicolumn{2}{c}{$+0.2\beta$}&\multicolumn{2}{c}{$-0.2\beta$}\\
$t$ (s)&$n(t)$&$c(t)$&$n(t)$&$c(t)$\\
\midrule
0 & $1.000000\times10^6$ & $1.625000\times10^9$ & $1.000000\times10^6$ & $1.625000\times10^9$\\
0.1 & $1.252012\times10^6$ & $1.627938\times10^9$ & $8.323657\times10^5$ & $1.622974\times10^9$\\
10 & $1.525874\times10^6$ & $1.984017\times10^9$ & $7.294773\times10^5$ & $1.422359\times10^9$\\
\bottomrule
\end{tabular}
\end{center}
Densities are in $\mathrm{cm^{-3}}$. Both densities are continuous at $t=0$.
The initial derivatives are $\dot n(0^+)=\rho_0 n_0/\Lambda$ and
$\dot c(0^+)=0$. Independent matrix-exponential and Radau calculations
verify the coefficients and sampled densities above.
\end{adaptedsolution}

\newpage

\noindent{\bf Problem 3}

Consider again the kinetics equations with a single precursor group (as
defined in Problem 2). Suppose a step reactivity \(\rho\) is inserted
into an initially critical reactor. Show that if

\[
\frac{1}{\lambda\Lambda}|\beta-\rho|\gg1
\quad\text{and}\quad
\frac{\beta}{\lambda\Lambda}\gg1,
\]

then

\[
n(t)=n(0)\left(
\frac{\rho}{\rho-\beta}e^{\frac{\rho-\beta}{\Lambda}t}
+\frac{\beta}{\beta-\rho}e^{\frac{\lambda\rho}{\beta-\rho}t}
\right).
\]

\begin{draftsolution}
\textbf{Source limitation and new derivation.} The source solution says only
``See me in office hours.'' The current student Problem 3 retains the
original hypotheses, which do not imply the claimed result universally.
The equality must be understood as an asymptotic approximation, with the
additional scale-separation restriction below. This is a new, explicitly
qualified derivation, not a recovered source answer. Instructor review of
the student prompt is still needed.

Put $d=\beta-\rho$. The exact equation from Problem 2 is
\[
\Lambda\omega^2+(d+\lambda\Lambda)\omega-\lambda\rho=0.
\]
Away from $d=0$, separated fast and slow scales require
\[
\frac{\lambda\Lambda}{|d|}\ll1,\qquad
\frac{\lambda\Lambda|\rho|}{d^2}\ll1.
\]
The second condition does not follow in general from the two conditions
printed in the source. Under these restrictions,
\[
\omega_f\simeq-\frac d\Lambda=\frac{\rho-\beta}\Lambda,
\qquad \omega_s\simeq\frac{\lambda\rho}d.
\]
Apply $a_f+a_s=n_0$ and
$\omega_f a_f+\omega_s a_s=\rho n_0/\Lambda$.
To leading order in the scale separation,
\[
a_f\simeq-\frac\rho d n_0=\frac\rho{\rho-\beta}n_0,
\qquad a_s\simeq\frac\beta d n_0,
\]
which proves the displayed approximation. It satisfies $n(0)=n_0$ exactly,
while its initial derivative agrees only to the retained asymptotic order.
For $\rho<\beta$, the fast mode decays; the amplitude of the remaining
slow mode is the prompt-jump estimate $n_0\beta/(\beta-\rho)$.

\textbf{Why the qualification matters.} For example, take
$\beta=0.0065$, $\lambda\Lambda=10^{-8}$, and $d=10^{-6}$.
Then $|d|/(\lambda\Lambda)=100$ and $\beta/(\lambda\Lambda)=650000$, yet
$\lambda\Lambda|\rho|/d^2=64.99$, so the proposed root separation fails.
As $\rho\to\beta$, use the exact quadratic; the displayed approximation
is singular there. This derivation replaces the legacy ``office hours'' answer.

For fixed $\beta>0$, set $\lambda\Lambda=\beta\epsilon$ and
$\delta=d=\beta\epsilon^{3/4}$. Then the stated ratios are
$1/\epsilon$ and $\epsilon^{-1/4}$, whereas
$4\lambda|\rho|\Lambda/\delta^2
=4(1-\epsilon^{3/4})/\sqrt{\epsilon}\to\infty$.
This is an asymptotic failure of the supplied assumptions, not merely a
choice of a borderline numerical meaning for $\gg1$.

\end{draftsolution}
